\chapter{Correctness and Security Analysis}
\label{chap:correctness-security}

\section{Functional correctness of the optimistic path}

This chapter analyzes implementation correctness and security at the level of the codebase
developed during the project. It is not a formal proof of the SOX protocol. The goal is more
practical: identify which properties are enforced by the implementation, which properties
are covered by tests, and which issues remain as limitations or future work.

The optimistic path is implemented as a state machine. The relevant states are payment,
key release, buyer-side dispute sponsor registration, vendor-side dispute sponsor
registration, dispute, and final state. Each transition checks both the current state and the
expected actor. In the direct normal ERC-4337-compatible contract, this check can be
satisfied either by the actor directly or through a validated UserOperation context. In the
specialized clone contracts, it is satisfied by direct calls from the expected address.

The precontract variants modify the optimistic path only where the protocol requires it.
The normal path starts by waiting for buyer payment. The no-deposit path also waits for
buyer payment, but removes the optimistic sponsor deposit and reduces the required buyer
payment to the agreed price. The \(S=B\) path fuses creation and buyer payment: the buyer is
the deployer, the buyer deposit is recorded during initialization, and the contract starts
waiting for the key. The \(S=V\) path requires the vendor to be the sponsor, but still waits
for the buyer payment.

Correctness of these transitions is checked at three levels. First, the frontend validates
obvious role constraints before deployment. Second, the factory and specialized contracts
repeat the same constraints on-chain. Third, the Hardhat tests verify the resulting state
and deposits. In particular, the tests check that:

\begin{itemize}
  \item the no-deposit variant records no sponsor deposit and rejects unexpected value;
  \item the \(S=B\) variant starts after the buyer payment has already been recorded;
  \item the \(S=V\) variant enforces the vendor as deployer and still waits for buyer
  payment;
  \item unsafe combinations, underfunded deployments, and unauthorized actor calls revert;
  \item the 12 role combinations studied in Phase 2 can be executed in the test matrix.
\end{itemize}

The application also contains a native precontract verification route. When the buyer views
a non-accepted precontract, the UI can call \path{/api/precontracts/verify}, which invokes
\path{check_precontract_cli}. This recomputes the ciphertext and circuit commitments from
the stored ciphertext, description, commitment, and opening value. Thus, the implementation
contains a concrete verification mechanism for the buyer-facing precontract data. The
current proof-of-concept interface should still be read as a local demonstration client, not
as a fully hardened wallet-integrated client.

\section{Functional correctness of the dispute path}

The dispute path begins only after the optimistic contract has reached the appropriate
state. The dispute constructor checks that the optimistic contract is waiting for the
vendor-side dispute sponsor, that a buyer-side dispute sponsor has already been set, that
the circuit version is supported, that the number of gates is non-zero, and that sufficient
funds are available. This prevents a dispute account from being initialized from an
inconsistent optimistic state.

The normal dispute contract implements the binary-search structure of Step 7. The interval
is represented by \(a\), \(b\), and \texttt{chall}. The buyer publishes a response for the
current challenge, the vendor gives an opinion, and the interval is updated until a boundary
or a single gate is reached. Depending on the final challenge value, the contract enters one
of the Step 8 states: left boundary, right boundary, or ordinary gate verification.

The Step 8 verification path checks several commitments at once. For an ordinary gate, the
contract opens the global commitment to recover \(h_{\mathit{circuit}}\) and
\(h_{\mathit{ct}}\). It checks that the supplied gate belongs to the circuit accumulator,
that ciphertext inputs used by the gate belong to \(h_{\mathit{ct}}\), that previous values
belong to the previous \(hpre\), and that the gate output extends the accumulator to the
current value. The Solidity implementation uses the V2 evaluator to recompute the selected
gate result on-chain.

The Step 9 logic is also tested in both normal and self-sponsored configurations. In the
normal dispute contract, Step 9 may relaunch the challenge loop when the losing party is not
the corresponding sponsor. In the self-sponsored dispute contract, this loop is specialized:
if \(S_B=B\) and \(S_V=V\), the contract moves directly toward completion or cancellation
after the Step 8 result. This is intentionally a small override of Step 9 rather than a full
rewrite of the dispute verifier.

Functional tests and performance tests both exercise the dispute path. The tests include
normal disputes, self-sponsored disputes, hardcoded disputes, same-scope comparisons, and a
native 1 GiB hardcoded benchmark that executes precontract generation, every challenge
round, Step 8 left-boundary submission, and finalization. These tests do not constitute a
formal proof, but they are strong regression tests for the implemented state transitions and
the measured scenarios.

\section{Correctness of the hardcoded SHA-256 mode}

The hardcoded \(\mathsf{SHA256}\) mode has a narrower correctness objective than the generic
circuit mode. It does not try to support arbitrary descriptions. It supports the case where
the description is \(desc=\mathsf{SHA256}(x)\), and where the circuit layout is completely
determined by public metadata.

The optimistic hardcoded contract fixes this mode at deployment time. It stores the
description hash, plaintext length, and ciphertext IV, and verifies that the supplied number
of blocks and number of gates match the hardcoded layout. This prevents the vendor from
waiting until Step 8 to choose between the generic and hardcoded interpretations.

The dispute contract then reconstructs the expected gate bytes through
\path{HardcodedSha256CircuitLib.sol}. The library determines whether the gate is an AES-CTR
block gate, a padding gate, a SHA-256 compression gate, the description constant, or the
final comparison gate. The strict hardcoded tests call Step 8 with empty generic circuit
inputs: \texttt{gateBytes = 0x} and an empty \(\pi_1\). The contract therefore does not rely
on a vendor-supplied circuit gate as the source of truth.

This does not mean that Step 8 becomes witness-free. The contract does not store the entire
ciphertext or all intermediate values. The vendor still supplies values needed to evaluate
the selected gate. These values are checked against \(h_{\mathit{ct}}\), previous \(hpre\),
and extension proofs before the result is accepted. The hardcoded improvement is therefore
specific: the contract determines \(g_i\), while the supplied \(v_i\)-style values remain
authenticated witnesses.

The hardcoded mode also reduces exposure to the generic circuit-encoding ambiguities
discussed below, because the gate bytes are reconstructed by the contract for this
particular SHA-256 layout. It does not, however, fix the generic V2 circuit format for other
descriptions.

\section{Sponsor safety and incentive preservation}

The project changed the sponsor structure, so it was important not to introduce a path where
one actor can force another actor into a disadvantageous state. The implementation addresses
this through on-chain role checks and through explicit authorization when a third party acts
as buyer-side dispute sponsor.

For precontract-time sponsor equalities, the checks are simple. The \(S=B\) factory path and
contract require the deployer to be the buyer. The \(S=V\) path requires the deployer to be
the vendor. The no-deposit path requires zero initial value and removes the sponsor
reimbursement accounting from the optimistic phase. These checks are enforced on-chain, not
only in the UI.

For dispute-time sponsorship, the delicate case is an external \(S_B\). The supervisor
correctly pointed out that an external buyer-side dispute sponsor must not be able to claim
that the buyer is unhappy without the buyer's consent. The implementation therefore includes
an authorization-based function for external \(S_B\). The signed statement binds the buyer's
authorization to the chain id, the optimistic contract address, the buyer address, the
chosen \(S_B\), and the commitment. This prevents reuse across different exchanges or
different proposed sponsors.

For \(S_B=B\) and \(S_V=V\), the implementation uses explicit self-sponsored entry points.
This makes the Step 4--5 fusion and the Step 9 specialization transparent. The gas benefit
therefore comes from removing unnecessary repeated dispute iterations, not from skipping
authorization checks or weakening the dispute verifier.

The economic incentive analysis remains within the scope of the implemented protocol model.
The code preserves the relevant deposits and payouts of the tested scenarios, but it does
not attempt to solve all off-chain availability or censorship risks. Those assumptions are
separated below.

\section{Encoding ambiguities in the inherited circuit format}

The V2 circuit format encodes each gate on exactly 64 bytes:

\[
  \text{opcode} \parallel \text{sons} \parallel \text{parameters}
  \parallel \text{zero padding}.
\]

Sons are signed 48-bit indices encoded on 6 bytes. Positive indices point to previous gates;
negative indices point to ciphertext input blocks. The generator uses opcodes for AES-CTR,
SHA2, CONST, XOR, and COMP. The Solidity evaluator checks the opcode, parameter length,
padding, arity-dependent constraints, and rejects zero gate indices in the extraction logic.

Nevertheless, the inherited format is not as clean as it should be for a final protocol
specification. The arity is not explicitly encoded. For some calls, the Solidity parser
infers arity from the position at which the remaining bytes become zero padding. For other
Step 8 calls, the expected arity is effectively determined by the number of supplied witness
values. This works for the generated circuits used in the tests, but it is not a robust
canonical encoding rule.

The problematic cases are exactly the cases identified during the project discussion. A
SHA2 gate may have one or two sons, and a CONST gate may have zero or one son. Without an
explicit arity or distinct opcodes, a byte string can be interpreted differently depending
on whether the parser treats a field as a son or as part of the constant payload. The current
padding checks reduce accidental ambiguity, but they do not give the format the clean
separation one would want in a paper-ready specification.

The clean fix is to define a new circuit version with unambiguous opcodes or an explicit
arity byte. For example, one could separate CONST-without-son from CONST-with-son, and
separate the first SHA-256 compression gate from subsequent SHA-256 compression gates. Such
a change would require updating the Rust circuit generator, the Solidity evaluator, the
hardcoded layout, tests, and all commitment computations. For this semester project, the
format was analyzed and documented, but the generic V2 encoding was not replaced because
that would have changed the baseline being measured.

\section{Merkle tree domain separation issues}

The second inherited cryptographic issue is the accumulator construction. The implementation
builds Merkle-style accumulators with \(\mathrm{keccak256}\). Leaves are hashes of encoded
objects, and internal nodes are hashes of two child hashes. In the current format, however,
leaf hashing and node hashing are not domain-separated, and an odd element at a layer is
copied upward unchanged.

This creates two well-known sources of ambiguity. First, copying a lonely node upward means
that two different tree shapes can lead to the same value if the proof format does not
fully determine the shape. Second, using the same hash domain for leaves and internal nodes
allows a 64-byte encoded leaf to have the same byte length as two 32-byte child hashes
concatenated. In a cryptographic specification, this is exactly the kind of ambiguity that
should be avoided by construction rather than argued away case by case.

The source code itself already contains a warning about the lonely-node case in the Rust
accumulator implementation. The Solidity verifier mirrors the same accumulator convention
for compatibility with the generated proofs. Therefore, this is not a bug introduced by the
new sponsoring or hardcoded variants. It is an inherited accumulator-design limitation.

The appropriate fix is a new accumulator version with explicit domain separation. A clean
design would use distinct prefixes, for example

\[
  \mathsf{leaf} = H(0x00 \parallel \mathsf{encodedLeaf}), \qquad
  \mathsf{node} = H(0x01 \parallel \mathsf{left} \parallel \mathsf{right}),
\]

and a separate rule for single-child promotion, such as hashing it with its own prefix
rather than copying it unchanged. This would need to be applied consistently to
\(h_{\mathit{circuit}}\), \(h_{\mathit{ct}}\), \(hpre\), proof generation, and Solidity
verification. It was not patched in the final implementation because it would create a new
commitment format and make the gas comparison with the initial implementation less direct.

\section{ERC-4337 trust and liveness assumptions}

The original implementation uses ERC-4337-style UserOperations and an Alto bundler in the
normal path. The contracts validate signatures through \texttt{validateUserOp}, store the
role associated with the last validated operation, and then allow the account to execute the
corresponding role action. This protects against arbitrary execution through the EntryPoint:
the recovered signer must match the buyer, vendor, or dispute sponsor role expected by the
state transition.

However, ERC-4337 also introduces liveness assumptions that are not part of the pure SOX
protocol. A bundler may be unavailable, may refuse to include a UserOperation, or may require
different gas policies. A paymaster or EntryPoint deposit can also become an operational
dependency. These are not cryptographic breaks of the smart contract state machine, but
they are relevant for a practical fair-exchange system.

For this reason, the final architecture separates the optimized clone variants from the
normal account-abstraction path. The clone variants do not claim ERC-4337 support; they are
used for direct-transaction gas measurements and for simplified modes. The normal path
remains available for demonstrating the inherited account-abstraction integration. This
separation avoids hiding bundler assumptions inside the gas-optimized variants.

\section{Remaining security limitations}

The implementation should be read as a research prototype, not as production software. The
following limitations remain after the project:

\begin{itemize}
  \item The generic V2 circuit encoding is not fully canonical because arity is not encoded
  explicitly for all gate families.
  \item The Merkle accumulators lack domain separation and use copy-up behavior for lonely
  nodes.
  \item The proof-of-concept UI uses predefined development addresses and does not implement
  a production wallet authentication model.
  \item Some cryptographic material and intermediate commitments are stored in browser local
  storage during the demonstration workflow.
  \item The backend stores ciphertext artifacts centrally. Commitments make tampering
  detectable, but the backend is still an availability dependency.
  \item The hardcoded mode is implemented only for the \(desc=\mathsf{SHA256}\) circuit
  layout, not for arbitrary descriptions.
  \item The 1 GiB hardcoded dispute benchmark demonstrates a specific full path and witness
  generation strategy; it is not a complete generic dispute UI for every possible large-file
  scenario.
  \item The work relies on extensive automated tests and measurements, but it does not
  include a machine-checked formal proof of the Solidity implementation.
\end{itemize}

These limitations do not invalidate the results of the project. They define the boundary of
the contribution. The implemented work improves the gas architecture, adds and tests the
main protocol variants, makes large-file experiments practical, and identifies the inherited
cryptographic encoding issues that should be addressed in a future protocol version.
